% ============================================================
% resume AIAgent EN.tex — English version, same structure as resume AIAgent.tex
% ============================================================

\documentclass[10pt,a4paper]{article}

% --- Load style & templates ---
\input{templates/style-aiagent-en.sty}
\input{templates/cv.latex}

% --- Local extensions: 4-col / 5-col layouts ---
\let\cvheadingThree\cvheading% save 3-col version for project sections
% Education: 4-column (Date / School / Degree / Major)
\newcommand{\cvheadingEdu}[4]{%
    \noindent
    \makebox[0.29\textwidth][l]{\textbf{#1}}% Date
    \makebox[0.36\textwidth][l]{\textbf{#2}}% School
    \makebox[0.12\textwidth][l]{\textbf{#3}}% Degree
    \makebox[0.23\textwidth][r]{\textbf{#4}}\par% Major
}
\newcommand{\cveventEdu}[5]{%
    \vspace{0.3em}
    \cvheadingEdu{#1}{#2}{#3}{#4}
    \ifx\relax#5\relax\else
        \vspace{0.06em}
        #5
    \fi
}
% Work Experience: 4-column (Date / Company / Location / Role+Lang)
\newcommand{\cvheadingWork}[4]{%
    \noindent
    \makebox[0.33\textwidth][l]{\textbf{#1}}% Date
    \makebox[0.24\textwidth][l]{\textbf{#2}}% Company (logo+name)
    \makebox[0.12\textwidth][l]{\textbf{#3}}% Location
    \makebox[0.31\textwidth][r]{\textbf{#4}}\par% Role+Lang
}
\newcommand{\cveventWork}[5]{%
    \vspace{0.2em}
    \cvheadingWork{#1}{#2}{#3}{#4}
    \ifx\relax#5\relax\else
        \vspace{0.04em}
        #5
    \fi
}
% Tech stack icons next to section titles (PNG color icons)
% Priority: devicon(-original) → simple-icons(-simple) → manual(-custom)
\newcommand{\ticon}[1]{%
    \IfFileExists{assets/myIcons/#1-original.png}{%
        \raisebox{-0.08cm}{\includegraphics[height=0.42cm]{assets/myIcons/#1-original.png}}%
    }{\IfFileExists{assets/myIcons/#1-simple.png}{%
        \raisebox{-0.08cm}{\includegraphics[height=0.42cm]{assets/myIcons/#1-simple.png}}%
    }{%
        \raisebox{-0.08cm}{\includegraphics[height=0.42cm]{assets/myIcons/#1-custom.png}}%
    }}%
}
\newcommand{\seclogo}[1]{\hspace{0.12em}\ticon{#1}}

% ============================================================
% Personal Information
% ============================================================
\name{Zhiwen Zhao}
\lineA{%
    \faPhone\ 185-5613-5826 \enspace|\enspace
    \faEnvelope\ \href{mailto:cnjszzw@qq.com}{cnjszzw@qq.com} \enspace|\enspace
    \faWeixin\ cnjszzw%
}
\lineB{Hometown: Yangzhou, Jiangsu \enspace|\enspace Age: 26 (1999/10/27) \enspace|\enspace Gender: Male}
\lineC{Job Preference: Software Development}
\photo{assets/photo.jpg}

% ============================================================
% Document Body
% ============================================================
\begin{document}

\makecvheader

% ============================================================
% Education
% ============================================================
\vspace{-0.825cm}% shift first section up to use header whitespace
\cvsection{Education}

% Switch to 4-column layout (Date / School / Degree / Major)
\let\cvheading\cvheadingEdu
\let\cvevent\cveventEdu

\cvevent{2022/09 -- 2025/01}{%
    \raisebox{-0.12cm}{\includegraphics[height=0.50cm]{assets/school/Xijing_University_Emblem.png}}%
    \hspace{0.4em}Xijing University}{M.S.}{Computer Technology}{}

\cvevent{2018/09 -- 2022/06}{%
    \raisebox{-0.12cm}{\includegraphics[height=0.50cm]{assets/school/Shandong_Normal_University_seal.png}}%
    \hspace{0.4em}Shandong Normal University}{B.S.}{Broadcasting \& TV Directing}{}

% ============================================================
% Work Experience
% ============================================================
\cvsection{Work Experience}

% Switch to 5-column layout (Date / Company / Location / Industry / Role+Lang)
\let\cvheading\cvheadingWork
\let\cvevent\cveventWork

\cvevent{2025/04 -- 2026/03 (Full-time)}{%
    \raisebox{-0.12cm}{\includegraphics[height=0.55cm]{assets/LOGO/上海环创通信科技股份有限公司.png}}%
    \hspace{0.4em}Huanchuang}{Shanghai}{(Java/Vue/Python) Full-Stack Dev}{}

\cvevent{2024/02 -- 2024/05 (Internship)}{%
    \raisebox{-0.12cm}{\includegraphics[height=0.55cm]{assets/LOGO/美团.png}}%
    \hspace{0.4em}Meituan}{Shanghai}{(Java) Test Dev}{}

\cvevent{2023/09 -- 2024/02 (Internship)}{%
    \raisebox{-0.12cm}{\includegraphics[height=0.55cm]{assets/LOGO/哔哩哔哩.png}}%
    \hspace{0.4em}Bilibili}{Shanghai}{(Python/Go) Test Dev}{}

% ============================================================
% Project Experience (grouped by company)
% ============================================================
\let\cvheading\cvheadingThree% restore 3-col version

\cvsection{Shanghai Huanchuang Communication Co., Ltd.}[%
    \seclogo{claudecode}\seclogo{cursor}%                                  AI Tools
    \seclogo{python}\seclogo{langchain}\seclogo{langgraph}\seclogo{deepseek}\seclogo{mcp}\seclogo{fastapi}% Python/AI Stack
    \seclogo{java}\seclogo{spring}\seclogo{springboot}\seclogo{mybatis}\seclogo{mybatis-plus}\seclogo{mysql}\seclogo{redis}\seclogo{mqtt}% Java Stack
    \seclogo{vuejs}\seclogo{javascript}\seclogo{typescript}\seclogo{cesium}% Frontend
    \seclogo{git}\seclogo{docker}\seclogo{jenkins}%    Infrastructure
]

\cvproject{Project 1: Drone Management Cloud Platform}{}{2025/04 -- 2026/03}{%
    \cvtech{Java, SpringBoot, Redis, MySQL, MyBatis(Plus), WebSocket, MQTT, Vue, JavaScript, TypeScript, Cesium}
    \cvlabeled{Overview}{Built on top of the company's existing IoT platform, leveraging the self-organizing mesh network's ultra-long-distance transmission (>20km) advantage to add a drone access module, achieving integrated air-ground command and dispatch.}
    \cvlabeled{Implementation}{%
        \begin{itemize}
            \item \textbf{Drone Access \& Device Management}
              \begin{starlist}
                \item[(S)Background:] Drones connect to base stations via self-organizing mesh hardware and need to be integrated into the existing IoT platform as a new device type. Unlike ground devices, their connection chain is longer and disconnection scenarios are more complex, so existing online/offline processes cannot be directly adapted.
                \item[(T)Task:] Integrate drones into the existing platform's device management system, reuse permission control and online/offline processes, and solve the problem of devices failing to auto-recover after server restart.
                \item[(A)Action:] Backend integrated drones into the existing device management flow, managed online status via Redis, distinguishing between active disconnection and abnormal disconnection. Leveraged Spring startup-ready event listeners to broadcast custom MQTT reconnect commands after service startup, triggering drones to re-initiate the online process.
                \item[(R)Result:] Drones fully integrated into the existing device management system, sharing permission control and online/offline logic with ground devices; automatic online recovery after service restart, no manual intervention required.
              \end{starlist}

            \item \textbf{Flight Control \& Camera System}
              \begin{starlist}
                \item[(S)Background:] Drone flight and camera commands span many types (covering takeoff, fly-to-point, orbit flight, return home, emergency stop, etc. for flight; and zoom, photo, recording, gimbal reset, frame zoom, AIM targeting, etc. for camera). Initially, for rapid validation, most commands were sent directly from frontend to hardware via MQTT—fast to develop but logic scattered across frontend components, hard to maintain and secure.
                \item[(T)Task:] Implement all flight and camera command functionality, progressively migrate the command channel from frontend-direct to backend-unified processing, and add flight safety validation to prevent misoperation.
                \item[(A)Action:] Implemented flight commands (takeoff, fly-to-point, orbit, waypoint flight, return home, emergency stop) and camera payload commands (zoom, photo, recording, gimbal reset, frame zoom, AIM targeting). Initially used frontend-direct MQTT for rapid validation, later migrated command logic to backend for unified processing. Backend added flight safety validation: altitude limit comparison, landing confirmation and cancellation mechanisms.
                \item[(R)Result:] Full coverage of flight and camera control capabilities; command channel uniformly consolidated to backend; flight safety validation effectively prevents risks such as over-altitude flight.
              \end{starlist}

            \item \textbf{Drone Live Streaming}
              \begin{starlist}
                \item[(S)Background:] Hardware streaming is unstable; video streams may be delayed or even unreachable. Multiple commanders need to simultaneously view the same drone feed.
                \item[(T)Task:] Under unreliable hardware conditions, achieve reliable multi-user live viewing while avoiding bandwidth waste when no one is watching.
                \item[(A)Action:] Implemented async confirmation mechanism based on Spring DeferredResult—upon user viewing request, immediately release the thread; backend asynchronously polls ZLM waiting for stream readiness, then returns playback URL via callback once confirmed. Also designed Redis reference counting strategy: maintain a set of viewing users per stream; decrement count on user disconnect; auto-stop pull stream when count reaches zero.
                \item[(R)Result:] Eliminated thread blocking caused by hardware latency; multi-user viewing without interference; bandwidth auto-reclaimed when no viewers.
              \end{starlist}
        \end{itemize}
    }
}

\cvproject{Project 2: WVP Copilot --- LLM-Based Drone AI Agent}{}{2025/11 -- 2026/03}{%
    \cvtech{Python, LangChain, LangGraph, SSE, MCP, FastAPI, DeepSeek API, Redis}
    \cvlabeled{Overview}{Designed for fire emergency command scenarios, WVP Copilot deeply integrates hardware/software capabilities with AI algorithms to create a drone intelligent agent, dramatically lowering the barrier to use. Users simply describe their needs in natural language (e.g., "Fly to Lujiazui, zoom to 7x, take a photo, return home"), and Copilot accurately understands the intent, automatically decomposes it into a sequence of drone commands, orchestrates execution, and completes remote control and mission dispatch through conversational interaction. Built with a two-tier architecture: Python (Agent orchestration layer) + Java (device control layer), decoupled via REST communication.}
    \cvlabeled{Implementation}{%
        \begin{itemize}
            \item \textbf{State Synchronization \& Deterministic Fault Recovery for Async Operations}
              \begin{starlist}
                \item[(S)Background:] LLM inference is synchronous (seconds), while drone execution is asynchronous (minutes). Tool functions must wait for flight execution to complete after issuing commands, but wait time varies with flight distance and communication latency, creating uncertainty.
                \item[(T)Task:] Establish reliable state synchronization between LLM's synchronous reasoning and drone's async execution; after timeout, deterministically judge whether flight succeeded to avoid downstream command errors from misjudgment.
                \item[(A)Action:] Tool functions internally encapsulate Redis polling waits, dynamically estimating timeout thresholds based on flight distance. On timeout, instead of immediately declaring failure, query the drone's actual OSD position for deterministic fallback—distance to target < 5m \thickarrow\ classify as arrived (state lag due to communication delay); position unchanged \thickarrow\ classify as hardware fault; position still changing \thickarrow\ extend wait. Intermediate states pushed only to frontend; LLM does not perceive them.
                \item[(R)Result:] Async wait uncertainty eliminated; timeout scenarios resolved deterministically via OSD position; LLM always makes decisions based on accurate state.
              \end{starlist}

            \item \textbf{Cascading Failure Prevention \& Safety Isolation for Multi-Step Task Chains}
              \begin{starlist}
                \item[(S)Background:] User commands may decompose into multi-step task chains (fly to location \thickarrow\ take photo \thickarrow\ return home). If any step fails without blocking, subsequent steps may execute under abnormal conditions, creating safety risks. Emergency stop and other safety operations must be completely isolated from LLM to guarantee absolute reliability.
                \item[(T)Task:] Implement cascading failure blocking for multi-step chains—stop all subsequent execution immediately on any step failure; require human confirmation for high-risk operations; isolate emergency stop from Agent and LLM.
                \item[(A)Action:] In the Plan-Confirm-Execute three-stage workflow, any step returning failure or triggering emergency stop immediately breaks and blocks subsequent execution. High-risk operations (takeoff, return home) require mandatory Human-in-the-Loop confirmation before execution. Emergency stop commands go through an independent MQTT hardware channel, bypassing Agent and LLM entirely; frontend simultaneously closes SSE connection for session-level isolation.
                \item[(R)Result:] Multi-step task chains have cascading failure blocking; high-risk operations have human fallback; emergency stop physically isolated from LLM for absolute reliability.
              \end{starlist}

            \item \textbf{State-Machine-Driven Sub-Task Orchestration Replacing LLM Temporal Decision-Making}
              \begin{starlist}
                \item[(S)Background:] Orbit flight is inherently a state machine (enter orbit \thickarrow\ capture photos at angles \thickarrow\ exit at $360^\circ$), but LLM inference latency (1--3s per call) at drone flight speed of 3m/s would cause photo capture point drift of 3--9 meters. ReAct loop's dependency on LLM for temporal judgment is physically infeasible in real-time control scenarios.
                \item[(T)Task:] Transfer temporal decision-making for real-time control tasks like orbit flight from LLM to deterministic state machines, while preserving LLM's intent understanding capability.
                \item[(A)Action:] Introduced LangGraph to model explicit conditional subgraphs—enter\_orbit \thickarrow\ orbit\_loop (conditional edge: $\ge 360^\circ$ \thickarrow\ exit) \thickarrow\ exit\_orbit—where execution cycle is governed by graph edges; LLM only responsible for intent routing, not temporal decisions. Main task chain (linear tasks) retains AgentExecutor to avoid over-engineering.
                \item[(R)Result:] Eliminated photo capture point drift caused by LLM inference latency; orbit timing precision guaranteed by graph state transitions; LLM's role narrowed to intent understanding.
              \end{starlist}

            \item \textbf{MCP External Tool Chain \& Composite Command Encapsulation}
              \begin{starlist}
                \item[(S)Background:] Drone control involves multiple external systems (Amap geocoding, device control APIs), and a single user intent (e.g., "record for 30 seconds") often maps to multiple low-level command sequences (start \thickarrow\ countdown \thickarrow\ stop). LLM directly orchestrating low-level commands is error-prone and inefficient.
                \item[(T)Task:] Integrate external tool chains, encapsulate composite commands to reduce LLM orchestration complexity, enabling LLM to program against business semantics rather than hardware primitives.
                \item[(A)Action:] Integrated Amap Server via MCP protocol with 15 map tools auto-registered; geocoding returns multiple candidates for LLM to arbitrate the best match based on context. Defined 16 drone tool functions using LangChain @tool decorator, following the principle of "exposing business semantics, not hardware primitives"—encapsulating composite commands (e.g., record\_for\_duration internally orchestrates start \thickarrow\ countdown \thickarrow\ stop as a complete sequence). SSE streaming push enables real-time observability of the tool invocation chain.
                \item[(R)Result:] LLM invokes tools at business-semantic level; composite commands simplify common operations from multi-step to single-step; SSE streaming output makes the operation process transparent and observable.
              \end{starlist}
        \end{itemize}
    }
}

% ============================================================
% Personal Projects
% ============================================================
\cvsection{Personal Projects}[\seclogo{claudecode}\seclogo{java}\seclogo{spring}\seclogo{springboot}\seclogo{mybatis}\seclogo{mysql}\seclogo{redis}\seclogo{git}\seclogo{elasticsearch}\seclogo{rocketmq}\seclogo{nginx}]

\cvproject{Project 1: Bilibili-Style Danmaku Video Website}{}{2024/05 -- 2026/05}{%
    \cvlabeled{Repository}{%
        \begingroup
        \hypersetup{urlcolor=primaryblue}%
        \url{https://github.com/Cnjszzw/imooc-bilibili} \textbar\ \url{https://github.com/Cnjszzw/bilibili-microservices}%
        \endgroup
    }
    \cvtech{SpringCloud Alibaba, Nacos, Gateway, OpenFeign, Seata, MyBatis, Redis, RocketMQ, Elasticsearch, WebSocket}
}

% ============================================================
% Professional Skills
% ============================================================
\cvsection{Professional Skills}

\begin{multicols}{2}
\cvskill{Java Ecosystem}{Proficient in Java, familiar with JVM memory model and concurrent programming; skilled in SpringBoot, MyBatis (Plus); knowledgeable in SpringCloud microservice stack (Nacos, OpenFeign, Gateway)}

\cvskill{Data Storage \& Retrieval}{Familiar with MySQL, with hands-on SQL tuning experience including N+1 query optimization (preloading + HashMap indexing), deferred join, and covering indexes; familiar with Redis caching strategies and data consistency assurance; familiar with Elasticsearch full-text search (Multi-Match, highlighting, multi-dimensional sorting)}

\cvskill{AI Agent Engineering}{Hands-on experience building LLM Agents from 0 to 1: proficient in LangChain, LangGraph, Function Calling, MCP protocol; skilled in Agent reliability design; proficient with Cursor/ClaudeCode}

\cvskill{Distributed Systems}{Knowledgeable in distributed locks (Redisson), distributed transactions (RocketMQ transactional messages / TCC / reliable eventual consistency), configuration centers (Nacos), and other common component design patterns and application scenarios}

\cvskill{Performance Optimization}{Hands-on core API optimization experience: reduced 1K+ Mesh device query latency from 7--8s to under 1s, reduced 300K-row log deep pagination from 5s to under 1s (5× performance gain); adept at identifying and resolving query and storage bottlenecks in large-data-volume scenarios}

\cvskill{Engineering Practices}{Familiar with Logback logging framework and Kibana troubleshooting, familiar with Linux + Docker basic operations; capable across the full lifecycle from requirements analysis to integration and delivery; experienced in cross-team collaboration across frontend, backend, and hardware/software}
\end{multicols}

\end{document}
